\begin{theorem}[nodal-cubic test family]\label{prop:nodal-counterexample}
\Refuted\ The p.~692 implication described after (5.2) is false as a formal
sheaf-theoretic statement, even for a proper family with smooth total space
and an irreducible reduced normal-crossings central fibre.
\end{theorem}

\begin{proof}
Let $\Delta$ be a sufficiently small disc with coordinate $t$ and put
\[
\mathcal C=
\{([x:y:z],t)\in\mathbb P^2\times\Delta:
y^2z=x^2(x+z)+tz^3\}.
\tag{5.8}
\]
This is proper over $\Delta$.  Its total space is smooth.  On $z\ne0$, the
partial derivative of the defining equation with respect to $t$ is
$-z^3\ne0$.  On $z=0$, the equation forces $x=0$; then $y\ne0$ and the
$z$-derivative is $y^2\ne0$.  The central fibre $C_0$ is the irreducible
nodal cubic $y^2z=x^2(x+z)$, with node $[0:0:1]$.

Set
\[
\mathcal X=\mathcal C\times\mathbb P^1,\qquad
S=C_0\times\mathbb P^1,
\]
and let $g:\mathcal X\to\Delta$ be projection.  Then $\mathcal X$ is a
smooth threefold, $g$ is proper, and $S$ is an irreducible reduced
normal-crossings surface.  Its normalization is
$\widetilde S\simeq\mathbb P^1\times\mathbb P^1$.

Let $p,q$ be the two inverse images of the node in the normalization of
$C_0$.  Fix $\lambda\in\mathbb C^*$ which is not a root of unity (for
example, $\lambda=2$).  Glue the fibres of
the trivial line bundle on the normalization at $p$ and $q$ by the scalar
$\lambda$.  This defines a line bundle
$A_\lambda\in\operatorname{Pic}^0(C_0)$ whose pullback is trivial.  It has
infinite order: its $n$th tensor power has gluing scalar $\lambda^n$, which
is never one.  It is not trivial: a global section pulls back to a constant
$c$ satisfying
$c=\lambda c$, so
\[
H^0(C_0,A_\lambda)=0.
\tag{5.9}
\]
Put $A=\operatorname{pr}_1^*A_\lambda$ on $S$.  Then
$\eta^*A\simeq\mathcal O_{\widetilde S}$ and $H^0(S,A)=0$.
Because $S$ is the full fibre,
$N^*_{S/\mathcal X}\simeq\mathcal O_S$, so the conormal-twist global
sections vanish.

The natural map
$\widetilde\Omega^1_S\to\eta_*\Omega^1_{\widetilde S}$ is injective:
its kernel is supported where normalization is not an isomorphism, while
the source is torsion-free by definition.  Projection formula and
$\eta^*A\simeq\mathcal O_{\widetilde S}$ therefore give
\[
H^0(S,\widetilde\Omega^1_S\otimes A)
\hookrightarrow
H^0(\widetilde S,\Omega^1_{\widetilde S})=0,
\tag{5.10}
\]
the last equality following from
$\widetilde S=\mathbb P^1\times\mathbb P^1$.
Nevertheless \cref{prop:conductor} produces
\[
0\ne s\in H^0(S,\Omega^1_{\mathcal X}|_S\otimes A).
\]
Thus both vanishings invoked in the printed reduction hold while the desired
ambient vanishing fails.  This test chooses the line bundle $A$ directly on
$S$; it does not assert that $A$ extends to a line bundle on the total space
$\mathcal X$.  It therefore refutes the claimed \emph{formal implication
between the three sheaf-cohomology groups} for a non-torsion $A$, but does
not by itself instantiate the narrower global-line-bundle situation on
p.~692.  The family (5.8) also does not satisfy all hypotheses of CDP
Theorem~2.2, so this is not a counterexample to that theorem.
\end{proof}

\subsection{Cartier-divisor duality without a normality hypothesis}

